\documentclass[10pt,a4paper]{article}
\usepackage[utf8]{inputenc}
\usepackage[english]{babel}
\usepackage{amsmath}
\usepackage{amsfonts}
\usepackage{amssymb}
\usepackage{graphicx}

\usepackage{glossaries}

\makeglossaries
\newacronym{vqe}{VQE}{Variational Quantum Eigensolver}

\usepackage[left=2cm,right=2cm,top=2cm,bottom=2cm]{geometry}
\author{Philipp Hanussek}
\title{Quantum Variational Algorithms}

\begin{document}
\maketitle

\section{Motivation}
Since the development of the quantum computers, many claims regarding quantum advantage have been made. Often, these advantage claims did not last long. In this work, we focus exclusively on the D-Wave Quantum Annealer, which leverages quantum technology to solve quadratically unconstrained binary optimization (QUBO) problems. D-Wave is among the first companies to commercially sell quantum computers. In 2025, researcher from D-Wave claimed to have solved a scientifically relevant problem fast using quantum technology, than it would have been possible on classical devices. However, it is currently an open question whether the D-Wave Quantum Annealer is able to solve other problems well, and possible faster than convential computers. In this work, we assess the D-Wave's ability to sample from a probability distribution, to approximate the ground state of different statistical mechanical models using variational algorithms.  
\subsection{Previous Work}
The D-Wave's Quantum Annealers performance scaling for finding the ground state (corresponding to the minimum of the quadratic binary cost function) has been extensively studied, with classical methods exhibiting better scaling and solution times. This has been done on a wide range of problems, ranging from cryptographic instances to physics inspired models.
\section{Methods}
In this section we briefly outline the main computational and algorithmic results used to gather results.
 \subsection{Sampling Methods}
 We use classical and quantum state-of-the-art methods to sample from the target distribution.
 \paragraph{Quantum Annealing}
 Quantum Annealing is inspired by 
\section{Variational Algorithms}
Variational algorithms are a class of hybrid quantum-classical optimization methods. We focus on the \acrfull{vqe} which is designed to find the ground state of a given Hamiltonian. It consists of three parts: the \textit{cost function} to be minimized, the \textit{ansatz}, which in our case is a Boltzman machine, and a \textit{sampler}, which draws configurations from the distribution defined by the current ansatz.
 
\end{document}	
